\scanpage{98}
Plus généralement, définissons
\[
\underline{\operatorname{Hom}}_{S/k}(F,G)=
\underline H^0_{S/k}(\underline{\operatorname{Hom}}_S(\unsure{F|G},G))
\]
\deleted{\unsure{Alors}} et \inserted{notons} les foncteurs dérivés de
$G\longmapsto\underline{\operatorname{Hom}}_{S/k}(F,G)$ par
\[
\underline{\operatorname{Ext}}^i_{S/k}(F,G).
\]
On a alors, si on admet que pour $G$ injectif,
$\underline{\operatorname{Hom}}_S(F,G)$ est acyclique pour
$\underline H^0_{S/k}$, \deleted{et pour $H^0_S$}, les suites spectrales
\margin{vérifier\figure{\begin{tikzpicture}\draw(0,0)--(0,1.1);\end{tikzpicture}}}
\[
\underline{\operatorname{Ext}}^{*\deleted{\unsure{p+q}}}_{S/k}(F,G)
\Longleftarrow\underline H^p_{S/k}
(\underline{\operatorname{Ext}}^q_S(\unsure{F|k},G))
\]
et compte tenu de
\[
\operatorname{Hom}_S(F,G)\simeq
H^0(k,\underline{\operatorname{Hom}}_{S/k}(F,G))
\]
et admettant que pour $G$ injectif,
$\underline{\operatorname{Hom}}_{S/k}(F,G)$ est acyclique pour $H^0(k,-)$,
on trouve
\margin{vérifier\figure{\begin{tikzpicture}\draw(0,0)--(0,1.1);\end{tikzpicture}}}
\[
\operatorname{Ext}^*(S;F,G)\Longleftarrow
H^p(k,\underline{\operatorname{Ext}}^q_{S/k}(F,G)).
\]
Enfin on aura
\[
\underline H^*_{S/k}(F)\simeq
\underline{\operatorname{Ext}}^*_{S/k}(\mathbb Z_S,F).
\]

Notations et remarques analogues en partant de faisceaux sur
$U\unsure{\subset|\in}S$, au lieu de $S$.
\mspage{97}

